\documentclass[../../main.tex]{subfiles}

\begin{document}

\section{序拓扑}\label{section:ch02-序拓扑}

若 $X$ 是一个全序集，那么可以应用序关系在 $X$ 上定义一个标准拓扑，称之为序拓扑。本节讨论序拓扑并且研究它的某些性质。

设 $X$ 是具有全序关系 $<$ 的一个集合，给定 $X$ 的两个元素 $a$ 和 $b$，$a < b$。则存在 $X$ 的 4 个子集，称为由 $a$ 和 $b$ 所决定的区间。它们是：
\begin{align*}
(a, b) = \{x \mid a < x < b\},
\end{align*}
\begin{align*}
(a, b] = \{x \mid a < x \leqslant b\},
\end{align*}
\begin{align*}
[a, b) = \{x \mid a \leqslant x < b\},
\end{align*}
\begin{align*}
[a, b] = \{x \mid a \leqslant x \leqslant b\}.
\end{align*}
这里采用的记号与你所熟悉的当 $X$ 是实直线时所用过的记号相似，只不过这里的区间是在任意全序集上的。第一种类型的集合称为 $X$ 的开区间（open interval），最后一种类型称为 $X$ 的闭区间（closed interval），第二种、第三种类型的集合则称为 $X$ 的半开区间（half-open interval）。用“开”字，是暗示在赋予 $X$ 某一个拓扑之后 $X$ 中的开区间都是开集。其他术语的用意类此。

\begin{definition}[序拓扑]\label{definition:ch02-序拓扑}
设 $X$ 是具有全序关系的一个集合，其元素多于一个。设 $\mathcal{B}$ 为下述类型所有集合的族：
\begin{enumerate}[(1)]
\item $X$ 的所有开区间 $(a, b)$.
\item 所有形如 $[a_{0}, b)$ 的区间，其中 $a_{0}$ 是 $X$ 的最小元（如果存在的话）.
\item 所有形如 $(a, b_{0}]$ 的区间，其中 $b_{0}$ 是 $X$ 的最大元（如果存在的话）.
\end{enumerate}
则族 $\mathcal{B}$ 是 $X$ 的某一个拓扑的基。此拓扑称为序拓扑（order topology）。

如果 $X$ 没有最小元，那就没有第二种类型的集合，如果 $X$ 没有最大元，那就没有第三种类型的集合。
\end{definition}

我们验证 $\mathcal{B}$ 满足作为一个基所要满足的条件。首先注意，对于 $X$ 的任意一个元素 $x$，它至少在 $\mathcal{B}$ 的一个元素中；最小元（如果存在）在所有类型 (2) 的集合中，最大元（如果存在）在所有类型 (3) 的集合中，其他的元素则在类型 (1) 的集合中。其次注意，上述各类型中任意两集合之交，仍然是属于上述类型的一种，或者是一个空集。其他需要验证的情形留给读者完成。

\begin{example}\label{example:ch02-标准序拓扑}
上一节定义的 $\mathbb{R}$ 的标准拓扑，恰好就是由 $\mathbb{R}$ 中的常用序关系给出的序拓扑。
\end{example}

\par
\begin{figure}[htbp]
\centering
\includegraphics[width=0.6\textwidth]{ch02-fig14-1}
\caption{$\mathbb{R} \times \mathbb{R}$ 字典序下的两种开区间}\label{figure:ch02-平面序拓扑}
\end{figure}

\begin{example}\label{example:ch02-平面字典序}
在字典序下研究集合 $\mathbb{R} \times \mathbb{R}$，其元素记作 $x \times y$（这样是为了避开记号上的困难）。$\mathbb{R} \times \mathbb{R}$ 既无最大元也无最小元，于是 $\mathbb{R} \times \mathbb{R}$ 上的序拓扑有一个以形如 $(a \times b, c \times d)$（其中 $a < c$ 或者 $a = c$，$b < d$）的所有开区间的族组成的基。这两种类型的区间画在\cref{figure:ch02-平面序拓扑}中。可以验证，仅由第二种类型的区间组成的子族也是 $\mathbb{R} \times \mathbb{R}$ 上序拓扑的一个基。请读者自己验证。

\end{example}

\begin{example}\label{example:ch02-Z离散序拓扑}
正整数集 $\mathbb{Z}_{+}$ 是具有最小元的全序集。$\mathbb{Z}_{+}$ 上的序拓扑是离散拓扑，因为它的每一个单点集是开集：若 $n > 1$，则单点集 $\{n\} = (n-1, n+1)$ 是一个基元素。若 $n = 1$，则单点集 $\{1\} = [1, 2)$ 也是一个基元素。
\end{example}

\begin{example}\label{example:ch02-两点序型}
具有最小元的全序集的另一个例子是在字典序下的集合 $X = \{1, 2\} \times \mathbb{Z}_{+}$。用 $a_{n}$ 记 $1 \times n$，用 $b_{n}$ 记 $2 \times n$，则 $X$ 可以表示成
\begin{align*}
a_{1}, a_{2}, \cdots; b_{1}, b_{2}, \cdots.
\end{align*}
$X$ 上的序拓扑不是离散拓扑。它的大部分单点集是开集，但是有一个例外，就是单点集 $\{b_{1}\}$。任何包含 $b_{1}$ 的开集，必定包含着包含 $b_{1}$ 的一个基元素（根据定义），而包含 $b_{1}$ 的任何一个基元素必含有序列 $a_{i}$ 中的点。
\end{example}

\begin{definition}[射线]\label{definition:ch02-射线}
若 $X$ 是一个全序集，$a$ 是 $X$ 的一个元素，则 $X$ 有 4 个子集称为由 $a$ 决定的射线（ray），它们是：
\begin{align*}
(a, +\infty) = \{x \mid x > a\};
\end{align*}
\begin{align*}
(-\infty, a) = \{x \mid x < a\};
\end{align*}
\begin{align*}
[a, +\infty) = \{x \mid x \geqslant a\};
\end{align*}
\begin{align*}
(-\infty, a] = \{x \mid x \leqslant a\}.
\end{align*}
前两种类型的集合称为开射线（open ray），后两种类型的集合称为闭射线（closed ray）。
\end{definition}

使用“开”这个词，是因为 $X$ 中的开射线在序拓扑下是开集。其余的情形类此。例如，我们来看射线 $(a, +\infty)$。若 $X$ 有最大元 $b_{0}$，则 $(a, +\infty)$ 等于基元素 $(a, b_{0}]$。若 $X$ 没有最大元，则 $(a, +\infty)$ 等于所有形如 $(a, x)$ 的基元素的并，其中 $x > a$。所以，不论在哪种情形下，$(a, +\infty)$ 总是开集。对于射线 $(-\infty, a)$ 来说也是如此。

所有开射线组成 $X$ 上序拓扑的一个子基。事实上，由于开射线在序拓扑下是开集，所以由开射线所生成的拓扑含于序拓扑。另一方面，序拓扑中的每一基元素为有限个开射线之交。而 $(a, b)$ 等于 $(-\infty, b)$ 与 $(a, +\infty)$ 之交，同时 $[a_{0}, b)$ 和 $(a, b_{0}]$（如果存在的话）都是开射线。因此，由开射线所生成的拓扑也包含了序拓扑。

\end{document}